\subsection{李代数的幂零根基}

定义 3。设 g 是李代数。g 的幂零根基是 g 的有限维简单表示的核的交。

注记。（1）设 \(\mathfrak{s}\) 是 \(\mathfrak{g}\) 的幂零根基。由于 \(\mathfrak{g}\) 的向量子空间的每个递降序列都稳定，所以存在有限个 \(\mathfrak{g}\) 的有限维简单表示，其核的交为 \(\mathfrak{s}\)。

这些表示的直和是半单的，且其核为 s。因此，g 的有限维半单表示的核所成的集合有最小元，即 s。

\begin{enumerate}
\def\labelenumi{(\arabic{enumi})}
\setcounter{enumi}{1}
\item
  根据 §4 第 3 号的命题 4 (c)，s 也是 g 的有限维表示的最大幂零理想的交。特别地，s 包含于 g 的最大幂零理想中，因而是 g 的幂零理想。
\item
  每个线性形式 \(\lambda\) 在 \(\mathfrak{g}\) 上且在 \(\mathcal{D}\mathfrak{g}\) 上为零；它都是以 \(\mathbf{K}\) 为表示空间的 \(\mathfrak{g}\) 的简单表示，从而 \(\lambda (\mathfrak{s}) = \{0\}\)。因此 \(\mathfrak{s} \subset \mathcal{D}\mathfrak{g}\)。另一方面，由注记 2，\(\mathfrak{s}\) 包含于 \(\mathfrak{r}\) 的根基 \(\mathfrak{g}\) 中。我们将证明 \(\mathfrak{s} = \mathfrak{r} \cap \mathcal{D}\mathfrak{g}\)。
\end{enumerate}

引理 1. 设 \(\mathbf{V}\) 是 \(\mathbf{K}\) 上的有限维向量空间，\(\mathfrak{g}\) 是 \(\mathfrak{gl}(\mathbf{V})\) 的子代数，使得 \(\mathbf{V}\) 是简单的 \(\mathfrak{g}\)-模，\(\mathfrak{a}\) 是 \(\mathfrak{g}\) 的交换理想。则 \(\mathfrak{a} \cap \mathcal{D}\mathfrak{g} = \{0\}\)。令 \(\mathbf{S}\) 为由 1 和 \(\mathcal{L}(\mathbf{V})\) 生成的 \(\mathfrak{a}\) 的子代数。

若 \(\mathfrak{b}\) 是 \(\mathfrak{g}\) 中包含于 \(\mathfrak{a}\) 的理想，且对所有 \(\operatorname{Tr}bs = 0\) 及所有 \(b\in \mathfrak{b}\) 有 \(s\in S\)，则特别地，根据 \(S\) 的定义，对每个整数 \(\operatorname{Tr}(b^n) = 0\) 有 \(n > 0\)，从而 \(b\) 是幂零的（《代数》，第 VII 章，§ 5，第 5 号，命题 13 的推论 4）；由于 \(\mathfrak{b}\) 的元素全都是幂零的，\(\mathfrak{b} = \{0\}\)（§ 4，第 3 号，引理 2）。我们首先将此应用于 \([\mathfrak{g},\mathfrak{a}]\) 的理想 \(\mathfrak{g}\)。若 \(x\in \mathfrak{g}\)、\(a\in \mathfrak{a}\)、\(s\in S\)，则 \(\operatorname{Tr}[x,a]s = \operatorname{Tr}(xas - axs) = \operatorname{Tr}x(as - sa) = 0\)，因为 \(as = sa\)；所以 \([\mathfrak{g},\mathfrak{a}] = \{0\}\)。因此 \(\mathfrak{g}\) 的元素与 \(\mathfrak{a}\) 的元素可交换，因而也与 \(S\) 的元素可交换。若 \(x,y\) 属于 \(\mathfrak{g}\) 且 \(s\in S\)，则

\[
\operatorname{Tr} [ x, y ] s = \operatorname{Tr} (x y s - y x s) = \operatorname{Tr} x (y s - s y) = 0
\]

由于 \(ys = sy\)，取 \(b\) 为理想 \(\mathcal{D}\mathfrak{g} \cap a\)，便得到 \(\mathcal{D}\mathfrak{g} \cap a = \{0\}\)。

定理 1. 设 \(\mathfrak{g}\) 是李代数，\(\mathfrak{r}\) 是其根基，\(\mathfrak{s}\) 是其幂零根基。则 \(\mathfrak{s} = \mathcal{D}\mathfrak{g} \cap \mathfrak{r}\)。

我们已经知道 \(\mathfrak{s} \subset \mathcal{D}\mathfrak{g} \cap \mathfrak{r}\)。因此只需证明：若 \(\rho\) 是 \(\mathfrak{g}\) 的有限维简单表示，则 \(\rho(\mathcal{D}\mathfrak{g} \cap \mathfrak{r}) = \{0\}\)。令 \(k\) 是最小的、\(\geqslant 0\) 且使 \(\rho(\mathcal{D}^{k}\mathfrak{r}) = \{0\}\) 的整数；记 \(\mathfrak{g}' = \rho(\mathfrak{g})\)、\(\mathfrak{a}' = \rho(\mathrm{D}^k\mathfrak{r})\)；由于 \(\mathcal{D}^k\mathfrak{r}\) 是 \(\mathfrak{g}\) 的理想，\(\mathfrak{a}'\) 是 \(\mathfrak{g}'\) 的理想；由于 \(\rho(\mathcal{D}^{k}\mathfrak{r}) = \{0\}\)，这个理想是交换的。若 \(\mathrm{V}\) 是 \(\rho\) 的表示空间，则 \(\mathfrak{g}' \subset \mathfrak{gl}(\mathrm{V})\)，且 \(\mathrm{V}\) 是简单的 \(\mathfrak{g}'\)-模。于是 \(\rho(\mathcal{D}\mathfrak{g} \cap \mathcal{D}^k\mathfrak{r}) \subset \mathcal{D}\mathfrak{g}' \cap \mathfrak{a}' = \{0\}\)。若 \(k > 0\)，则 \(\mathcal{D}^k\mathfrak{r} \subset \mathcal{D}\mathfrak{g}\) 且 \(\rho(\mathcal{D}^k\mathfrak{r}) = \{0\}\)，这与 k 的定义矛盾。因此 \(k\) 为 \(k = 0\)，即

\[
\rho (\mathcal{D}\mathfrak{g} \cap \mathfrak{r}) = \{0 \}.
\]

推论 1. 设 g 是可解李代数。g 的幂零根基是 \(\mathcal{D}\mathfrak{g}\)。若 \(\rho\) 是 g 的有限维简单表示，则 \(\rho(\mathfrak{g})\) 是交换的，并且由 1 和 \(\rho(\mathfrak{g})\) 生成的结合代数 L 是 K 的有限次扩域。

这里 \(\mathfrak{r} = \mathfrak{g}\)，从而 \(\mathfrak{s} = \mathcal{D}\mathfrak{g}\)。所以 \(\rho(\mathcal{D}\mathfrak{g}) = \{0\}\)，这说明 \(\mathfrak{g}' = \rho(\mathfrak{g})\) 是交换的。根据 Schur 引理，\(\neq 0\) 的每个 \(\mathbf{L}\) 元素都是可逆的；所以 \(\mathbf{L}\) 是域。

推论 2（李定理）. 设 \(\mathfrak{g}\) 是可解李代数；假设 \(\mathbf{K}\) 代数闭。令 \(\mathbf{M}\) 是在 \(\mathfrak{g}\) 上有限维的 \(\mathbf{K}\)-模，\((\mathrm{M}_1)_{0 \leqslant i \leqslant r}\) 是 \(\mathbf{M}\) 的 Jordan-Hölder 序列。则对 \(\mathrm{M}_{i - 1} / \mathrm{M}_i\)，\(\mathbf{K}\) 在 \(1 \leqslant i \leqslant r\) 上的维数为 1，并且对所有 \(x \in \mathfrak{g}\)，有 \(x_{\mathrm{M}_{i - 1} / \mathrm{M}_i} = \lambda_i(x)\) . 1，其中 \(\lambda_i\) 是 \(\mathfrak{g}\) 上在 \(\mathcal{D}\mathfrak{g}\) 上为零的线性形式。特别地，每个在 \(\mathfrak{g}\) 上有限维的简单 \(\mathbf{K}\)-模实际上都是一维的。

令 \(\rho_{i}\) 是 \(g\) 在 \(\mathbf{M}_{i-1}/\mathbf{M}_{i}\) 上的表示。由 1 和 \(\mathbf{L}_{i}\) 生成的结合代数 \(\rho_{i}(g)\) 是域，是 K 的有限扩域，因而等于 K；而 \(\mathbf{M}_{i-1}/\mathbf{M}_{i}\) 是简单的 \(\mathbf{L}_{i}\)-模，所以 \(\dim \mathbf{M}_{i-1}/\mathbf{M}_{i} = 1\)。推论的其余部分显然。

注记。（1）若 \((\mathrm{M}_i)_{0\leqslant i\leqslant r}\) 被 M 的另一个 Jordan-Hölder 序列替代，则序列 \((\lambda_1,\dots ,\lambda_r)\) 被形如 \((\lambda_{\pi (1)},\dots ,\lambda_{\pi (r)})\) 的序列替代，其中 \(\pi\) 是 \(\{1,\dots ,r\}\) 的一个置换；这是 Jordan-Hölder 定理的结果。

\begin{enumerate}
\def\labelenumi{(\arabic{enumi})}
\setcounter{enumi}{1}
\tightlist
\item
  设 \((e_1, \ldots, e_r)\) 是 M 的一组基，满足 \(e_i \in \mathbf{M}_{i-1}, e_i \notin \mathbf{M}_i\)、\(1 \leqslant i \leqslant r\)（\(x \in \mathfrak{g}\)）。若 \(x\)，则相应于 x 的 M 的自同态相对于此基由一个三角矩阵表示，其对角系数为
\end{enumerate}

\[
\lambda_ {1} (x), \dots , \lambda_ {\tau} (x).
\]

推论 3. 假设 K 代数闭。若 g 是 r 维可解李代数，则 g 的每个理想都是维数为 r、r - 1、\ldots、0 的理想递降序列中的一项。

每个理想都是 g（看作伴随表示的表示空间）的一个 Jordan-Hölder 序列的一部分（《代数》，第 I 章，§ 6，第 14 号，定理 8 的推论）；因此只需应用推论 2。

推论 4. 假设 \(\mathbf{K} = \mathbf{R}\)。设 \(\mathfrak{g}\) 是可解李代数。每个 \(\mathfrak{g}\) 的简单表示的维数都 \(\leqslant 2\)。\(\mathfrak{g}\) 的每个理想都是理想递降序列 \((\mathfrak{g}_i)_{0\leqslant i\leqslant m}\) 中的一项，其中 \(\mathfrak{g}_0 = \mathfrak{g},\mathfrak{g}_m = \{0\}\)，且 \(\dim \mathfrak{g}_{i - 1} / \mathfrak{g}_i\leqslant 2\)（\(1\leqslant i\leqslant m\)）。

证明方式与推论 2、3 的证明类似，只需使用每个 \(\mathbf{R}\) 的代数扩域的次数都 \(\leqslant 2\) 这一事实。

推论 5. 李代数 g 可解的充要条件是 \(\mathcal{D}\mathfrak{g}\) 幂零。

必要性由推论 1 得到。充分性成立是因为 \(\mathfrak{g} / \mathcal{D}\mathfrak{g}\) 是交换的。

推论 6. 设 \(\rho\) 是李代数 \(\mathfrak{g}\) 的有限维表示。设 \(\mathfrak{r}\) 是 \(\mathfrak{g}\) 的根基。每个满足 \(x \in \mathfrak{r}\) 幂零的 \(\rho(x)\) 都属于 \(\mathfrak{n}\) 的最大幂零理想 \(\rho\)。

设 \(V\) 是 \(\rho\) 的表示空间；令 \(\langle V_i \rangle_{0 \leqslant i \leqslant r}\) 是 \(r\)-模结构在 \(V\) 上的 Jordan-Hölder 序列，并令 \(\rho_i\) 是 \(r\) 在 \(V_i / V_{i-1}\) 上的表示（\(1 \leqslant i \leqslant r\)）。若 \(\rho(x)\) 是幂零的，则 \(\rho_i(x)\) 也是幂零的；因为对所有 \(i\)，由 \(\rho_i(x)\) 生成的代数是域，所以 \(\rho_i(x) = 0\)。反之，若 \(\rho_i(x) = 0\)，则对所有 \(i, \rho(x) = 0\) 都成立。这说明集合 \(a\) 中满足 \(x \in r\) 且 \(\rho(x)\) 幂零的元素是 \(r\) 的理想。另一方面，\([g, a] \subset \mathcal{D}g \cap r \subset n \cap r \subset a\)，所以 \(a\) 是 \(g\) 的理想。这证明了 \(a \subset n\)。

推论 7. 设 g 是李代数，r 是其根基。以下四个集合相同：(a) g 的最大幂零理想；(b) r 的最大幂零理想；(c) 满足 \(x \in \mathfrak{r}\) 幂零的 \(\operatorname{ad}_{\mathfrak{g}} x\) 的集合；(d) 满足 \(x \in \mathfrak{r}\) 幂零的 \(\operatorname{ad}_{\mathfrak{r}} x\) 的集合。

将这些集合记为 \(\mathfrak{a}, \mathfrak{b}, \mathfrak{c}, \mathfrak{d}\)。包含关系 \(\mathfrak{a} \subset \mathfrak{b} \subset \mathfrak{d} \subset \mathfrak{c}\) 显然。由推论 6 应用于 \(\mathfrak{c} \subset \mathfrak{a}\) 的伴随表示，得 \(\mathfrak{g}\)。

